%; whizzy chapter -dvi
% -initex iniptex -latex platex -format platex -bibtex jbibtex -fmt fmt
% 以上 whizzytex を使用する場合の設定。

%     Tokyo Debian Meeting resources
%     Copyright (C) 2012 Junichi Uekawa
%     Copyright (C) 2011, 2015, 2020 Nobuhiro Iwamatsu

%     This program is free software; you can redistribute it and/or modify
%     it under the terms of the GNU General Public License as published by
%     the Free Software Foundation; either version 2 of the License, or
%     (at your option) any later version.

%     This program is distributed in the hope that it will be useful,
%     but WITHOUT ANY WARRANTY; without even the implied warranty of
%     MERCHANTABILITY or FITNESS FOR A PARTICULAR PURPOSE.  See the
%     GNU General Public License for more details.

%     You should have received a copy of the GNU General Public License
%     along with this program; if not, write to the Free Software
%     Foundation, Inc., 51 Franklin St, Fifth Floor, Boston, MA  02110-1301 USA

%  preview (shell-command (concat "evince " (replace-regexp-in-string "tex$" "pdf"(buffer-file-name)) "&"))

%%ここからヘッダ開始。

\documentclass[mingoth,a4paper]{jsarticle}
\usepackage{monthlyreport}
% 日付を定義する、毎月変わります。
\newcommand{\debmtgyear}{2022}
\newcommand{\debmtgmonth}{3}
\newcommand{\debmtgdate}{19}
% started from zero:
% (let ((year 2021) (month 1)) (+ (* (- year 2005) 12) month -1))
% and add 1
\newcommand{\debmtgnumber}{207}

% Needed to import pandoc-generated LaTeX documents.
% See https://stackoverflow.com/questions/40438037/tightlist-error-using-pandoc-with-markdown
\providecommand{\tightlist}{%
  \setlength{\itemsep}{0pt}\setlength{\parskip}{0pt}}

\begin{document}

\begin{titlepage}
\thispagestyle{empty}
% タイトルページ:編集必要な部分は最初のマクロに飛ばすこと

\vspace*{-2cm}
第\debmtgnumber{}回 東京エリア Debian 勉強会資料\\
\hspace*{-2cm}
\includegraphics{image-assets/dotdeb.pdf}\\
\hfill{}\debmtgyear{}年\debmtgmonth{}月\debmtgdate{}日

% ここはアップデートすること
% 全角文字にしないとフォントのサイズが合わないので注意
\rotatebox{10}{\fontsize{30}{30} {\gt ｇｉｔ－ｂｕｉｌｄｐａｃｋａｇｅ特集}}\\

\vspace*{-2cm}
\hfill{}\includegraphics[height=6cm]{image-assets/openlogo-nd.eps}
\end{titlepage}

\newpage

\begin{minipage}[b]{0.2\hsize}
 \definecolor{titleback}{gray}{0.9}
 \colorbox{titleback}{\rotatebox{90}{\fontsize{80}{80} {\gt デビアン勉強会} }}
\end{minipage}
\begin{minipage}[b]{0.8\hsize}
\hrule
\vspace{2mm}
\hrule
\begin{multicols}{2}
\tableofcontents
\end{multicols}
\vspace{2mm}
\hrule
\end{minipage}

\dancersection{最近のDebian関連のミーティング報告}{杉本 典充}

\subsection{2022 年 2 月度 東京エリア・関西合同Debian勉強会}

2022 年 2 月 19 日 (土) に東京エリア Debian 勉強会と関西 Debian 勉強会の合同で
オンラインによる Debian 勉強会を開催しました。参加者は 20 名でした。

% セミナー
セミナーは林さんの「技術書典12でDebianの薄い本を頒布した話」、taiseiyo さんの
「Debian で始める Python Programming」の2本を発表しました。

勉強会の終了後、参加者同士で Debian や OSS に関する話の情報交換を行いました。


\dancersection{事前課題}{杉本 典充}

今回の事前課題は以下です。

\begin{enumerate}
  \item git-buildpackage を知っていますか。
  \item Debianに関連して知りたいこと、聞いてみたいことがあれば自由に書いてみてください。
\end{enumerate}

%この課題に対して提出いただいた内容は以下です。

\begin{multicols}{2}
{\small
\begin{prework}{ ����(yy\_y\_ja\_jp) }
���ޤ���������ǤϤʤ��Ǥ���... �긵�Υ�åץȥåפǤ� /boot �ˤ� ext2
 ����¾�ˤ� ext3���Ƕ�ȤäƤ���ǥ����ȥåפǤ� /boot �ˤ� ext3 ����¾
 �ˤ� LVM ��� ext4 ��ȤäƤ��ޤ���
\end{prework}

\begin{prework}{ �����ϥ� }
�ǥե���Ȥ�ext3����
�����ǥե���ȤΤޤޡ�
(NTFS���VM���᡼�����ext3�⤢�뤱��)
\end{prework}

\begin{prework}{ yos.takahashi }
ext3/4���˻ȤäƤޤ���ext3�Υǡ�������١����ˤĤ�������Linux2011ǯ1���˼�ɮ���ޤ�����
\end{prework}

\begin{prework}{ MATOHARA }
����inode �ϳ���������Ƥ���inode ��ưŪ�˳�����Ƥ���XFS �����򤹤뤳
 �Ȥ�¿���Ǥ���NILFS �Ͼ�����Ƥߤ��ΤǤ�����mount ���˰ʲ��Τ褦�ʥ��
 ���������ФƤޤ��ݤ��ʤȻפ��ޤ�����
\begin{commandline}
$ sudo mount /dev/sdb1 /mnt
mount.nilfs2: WARNING! - The NILFS on-disk format may change at any time.
mount.nilfs2: WARNING! - Do not place critical data on a NILFS filesystem. 
\end{commandline}
����¾NotePC �Ǥ�dm-crypt �ξ�˥ե����륷���ƥ���֤��ưŹ沽�����ꡢ
 eCryptfs �ǰŹ沽�����ꤷ�Ƥ��ޤ����񤭹��߻���CPU �򤫤ʤ���񤷤ޤ��ġ�
\end{prework}

\begin{prework}{ ��ޤ� }
��ext2-$>$reiserfs-$>$jfs-$>$xfs-$>$reiserfs-$>$ext3�ȻȤäƤ��ޤ�����

����:
\begin{itemize}
 \item reiserfs: ���ե������¿���ե�����Υ������������ӥ��Ӥ��Ƥ��ɤ���
       �����������դȤ����ⵤ�������θ�μ�žȬ�ݥ����ɤˡ�����
 \item jfs: ����v1.0��̾��ä�IBM���ꥨ�ʥ��������ˤ�xfs��ƨ��
 \item xfs: fsck==true�˴�ư����������ǯ�Ȥä���ΤΥޥ�����Ĵ����0byte
       �ե��������������Ѥ���줺ƨ����
\end{itemize}
������reiser4��Ķ���Ԥ��뤦���ˤ��줬�����ʤäơ����ext3�˸��경��htree�����ä������⤦Ŵ�Ĥʤ鲿�Ǥ⤤���Ǥ��������Ȥ����Ĥ�nilfs�ʤɤ˼��Ф��Ƥ��ޤ���ext3�������noatime���٤Ǥ����������aufs���碌�Ƽ�ʬ���Ȥ��Ȥ� *strap �Ķ��򥯥����˥󥰤��Ƥ��ޤ����¸����ƥ��Ȥ������Ǥ�����Ǥ���USB�����ư�Ǥ�ͭ�ѡ�

���LVM�ǤϤʤ�MD��Ȥäƾ�Ĺ���ܥХå����åפ򤷤Ƥ��ޤ���
 MD(sda,sdb,sdc)�ǹ��ۤ����̾��MD(����)�Dz�ư���Хå����åפλ���
 attach/detach�򤹤롣�֥��å���٥�ʤΤǥꥫ�Х��FSǤ���Ǥ��������̤�
 �������¾����ˡ���ʤ�������
\end{prework}

\begin{prework}{ henrich }
�ȤäƤ���֤�NTFS��Ĺ���󤸤�ʤ��Ǥ����͡����졣
�����Ρ����Ѥο������ǥ�������ext4�ǥե����ޥåȤ��ޤ����������㤤��Ƚ��ʤ��Ȥ����������Ƥ��ޤ���
\end{prework}

\begin{prework}{ emasaka }
�Ĥ뤷��FS��ȤäƤޤ�
\end{prework}

\begin{prework}{ �ܾ� }
ext3����Ѥ��Ƥ��ޤ����ä��Ѥ�ä����ȤϤ��Ƥ��ޤ���
FS����ʤ��Ǥ������Ƕ�Lenny��2TB��HDD��Ȥä���parted�äƤλȤ��ƶä��ޤ�����
\end{prework}

\begin{prework}{ ����@������ }
���������Ū�˳��Ѥ��Ƥ���ե����륷���ƥ��ReiserFS�Ǥ���
�Ż��ǻȤäƤ�Ķ���ext3�Ǥ�����ext3���Ÿ������ǥ��㡼�ʥ뤬
����Ʋ��Ǥ������ηи��ʸŤ������ͥ�Ǥ���...�ˤ����ꡢ
���ޤ꿮�Ѥ��Ƥ��ޤ���
�����ReiserFS�Ķ��ǤϤ���ޤǤνꤤ���ʤ��Ÿ����ڤä��ꡢ
�Ƥ�HDD�����줫�����ꤷ�Ƥ��ﳲ����ä��и���̵���Τǡ�
��³Ū�˻ȤäƤ��ޤ���
��ǯ����ReiserFS�Υᥤ��ȯ��(Hans Reiser)�����ᤵ��Ƥ��ޤ������ƥʥ󥹤��ۤ��Ƥ��ޤ�����
�����������θ��ReiserFS��¾�γ�ȯ�Ԥˤ���³�����ݼ餵��Ƥ���Τǡ��¿����ޤ�����
\end{prework}

\begin{prework}{ nozzy123nozzy }
\begin{enumerate}
 \item LVM�ˤĤ��Ƥϡ�CentOS5.5��Ƴ�������Τ��Τޤޤ����Ѥ��Ƥޤ�����
       ���������ƥब��äƤ���Volume̾�ϥǥե���Ȥ�����ȡʼºݤˤ�
       kickstart�ˤơˤ��ѹ����ƻȤäƤޤ����ʾ㳲���Υ���١����˺��뤿
       ���
\item ext3�ˤĤ��Ƥϡ�debian-sid�����Τޤ޻��ꤷ�Ƥ����Τ򤽤Τޤ޻Ȥ�
      �Ƥ����ꤷ�ޤ���relatime, noatime ���餤�Ͼ����ɲä��Ƥߤ����ʡ���
      �ϻפäƤޤ���
\end{enumerate}
\end{prework}

\begin{prework}{ �ޤ��������ؤ� }
\begin{itemize}
 \item Debian�Ǥ��ä˶Ťä����Ȥ�����ext3��ȤäƤޤ������ۥޥ����qcow2
       ���᡼���ǥ��������ѻ��ʳ��ϡ�LVM�ϻȤäƤޤ���
 \item �����ǰ����ü�ʤΤϡ��������DHCP�������Ѥ�Armadillo-J�ǻȤäƤ�
       ��JFFS�Ǥ����ǥե���ȤΥե����०�����Ǥϥ�֡��Ȥ�����������
       �ƽ��������Ƥ��ޤ��Τǡ�RAM�ΰ�˽񤭤��ߡ��Ÿ��ڤäƤ�ä��ʤ�
       ���������Ǥ���Debian��udhcp�Υ������ѥå���������ӥ�ɤ��ƻȤäƤޤ���
       \footnote{\url{http://d.hatena.ne.jp/mkouhei/20080601/1212330630}}
 \item ��Debian���ߤǡ���ʬ����ǰ��֥ۥåȤʤΤ�palm webOS�Ǥ��������Ubuntu��
       �������ޥ���������Τ餷���ΤǤ�����/etc/mtab�򸫤��35�Ԥ⤢�ꡢ
       ���ʤ����֤ʹ����Ǥ��͡�
\end{itemize}
\end{prework}
}
\end{multicols}

%\dancersection{Debian Trivia Quiz}{username}
%
%Debianの昨今の話題についてのQuizです。
%
%今回の出題範囲は\url{debian-devel-announce@lists.debian.org} や \url{debian-news@lists.debian.org}などに投稿された内容からです。
%
%\begin{multicols}{2}
%%; whizzy-master ../debianmeetingresume201211.tex
% $B0J>e$N@_Dj$r$7$F$$$k$?$a!"$3$N%U%!%$%k$G(B M-x whizzytex $B$9$k$H!"(Bwhizzytex$B$,MxMQ$G$-$^$9!#(B
%

\santaku
{DebConf13 $B$N3+:ECO$H3+:EF|$O!)(B}
{$BF|K\!"El5~ET(B 6$B7n(B20$BF|(B}
{$B%K%+%i%0%"(B $B%^%J%0%"(B 7$B7n(B8-14$BF|(B}
{$B%9%$%9!"%t%)!<%^%k%-%e(B 8$B7n(B11-18$BF|(B}
{3}
{$B%K%+%i%0%"$O(BDebConf12$B$N3+:ECO$G$9!#(B
DebConf13$B$O%9%$%9$N%-%c%s%WCO$G3+:E$G$9!#(B
6/20$B$O3'$5$sM=Dj$r6u$1$F$*$-$^$7$g$&!#(B}

\santaku
{$B@$3&$N(BWeb$B%5!<%P$G:G$b?M5$$N$"$k(BLinux $B%G%#%9%H%j%S%e!<%7%g%s(B(W3Techs$BD4$Y(B)$B$O!)(B}
{CentOS}
{Debian}
{Ubuntu}
{B}
{\url{http://w3techs.com/technologies/history_details/os-linux}$B$K7k2L$N%0%i%U$,$"$j$^$9!#(B
$B8=:_(B Linux $B$r;HMQ$7$F$$$k(B web $B%5!<%P$N(B 32.9\% $B$,(B Debian $B$rMxMQ$7$F$*$j!"$=$N3d9g$O8=:_$bA}2C$rB3$1$F$$$k$=$&$G$9!#(B}

\santaku
{Debian $B%+!<%M%k%A!<%`$N%a%s%P!<$G$"$j!"(Bkernel.org $B$N(B 3.2.y $B0BDjHG7ONs$N%a%s%F%J$G$b$"$k(B Ben Hutchings $B$5$s$,<!4|(B Debian $B0BDjHG$H0l=o$K=P2Y$5$l$k(B Linux $B%+!<%M%k$K(B (3.2 $B7ONs$N(B mainline $B$K$OL5$$(B) $BDI2C5!G=$,Ek:\$5$l$kM=Dj$G$"$k$H=R$Y$F$$$^$9!#(B
$BB?$/$NDI2CE@$NCf$K4^$^$l$J$$$b$N$O2?!)(B}
{PREEMPT\_RT}
{Hyper-V guest drivers$B$N6/2=(B}
{ARM64/AArch64$B%"!<%-%F%/%A%c%5%]!<%H(B}
{C}
{Hyper-V guest drivers$B$O(Bmainline kernel$B$G(B3.2$B$K$b4^$^$l$F$$$^$9$,!"$h$j2~A1$5$l$?(B3.4$B$+$i$N=$@5$,F3F~$5$l$^$9!#(B
PREEMPT\_RT$B$O%O!<%I%j%"%k%?%$%`$r<B8=$9$k$?$a$N(BPatch$B!"(B
linux-image-rt-amd64 , linux-image-rt-686-pae $B$N(Bmetapackage$B$G;HMQ$G$-$^$9!#(B
$B?7$7$$(BARM 64$B%S%C%H%"!<%-%F%/%A%c%5%]!<%H$O(Bmainline kernel 3.7$B$+$i(B}

\santaku
{Wookey$B$5$s$,%"%J%&%s%9$7$?(Balpha$BHG$N(BDebian port arm64 image$B$O!)(B}
{Debian/Ubuntu port image}
{Debian/KFreeBSD port image}
{Debian/GnuHurd port image}
{A}
{self-bootstrapp(non x86)$BBP1~$H$N$3$H$G$9!#(B\url{http://wiki.debian.org/Arm64Port}$B$G%9%F!<%?%9$,3NG'$G$-$^$9!#(B}

\santaku
{700,000$BHVL\$N%P%0$,Js9p$5$l$?F|$rEv$F$k(B700000thBugContest$B$N7k2L$,=P$^$7$?!#$=$NM=A[F|$HJs9pF|$O!)(B}
{2012/12/12$B$rM=A[$7$?(BDavidPrevot}
{$BM=A[F|(B:2013/02/04$B!"Js9pF|(B:2013/02/14}
{$BM=A[F|(B:2013/02/07$B!"Js9pF|(B:2013/02/14}
{$BM=A[F|(B:2013/02/14$B!"Js9pF|(B:2013/02/07}
{C}
{$B:G$b6a$$(B2013/02/14$B$rM=A[$7$?(BChristian Perrier$B$5$s$,Ev$F$^$7$?!#7k2L$O(B\url{http://wiki.debian.org/700000thBugContest}$B$G8x3+$5$l$F$$$^$9!#(B
$B$^$?!"(B800,000/1,000,000$BHVL\$N%P%0$,Js9p$5$l$kF|$rEv$F$k%3%s%F%9%H(B\url{http://wiki.debian.org/800000thBugContest}$B$b3+:E$5$l$F$$$^$9!#(B}

\santaku
{master.debian.org$B$,?7$7$$5!3#$K0\9T$5$l$^$7$?!#$3$l$O2?$N%5!<%P$G$7$g$&$+(B $B!)(B}
{@debian.org$B$N%a!<%k%5!<%P(B}
{$B%Q%C%1!<%8$N%^%9%?!<%5!<%P(B}
{$B%Q%C%1!<%8$N%9%]%s%5!<(B(mentor)$B$rC5$9%5!<%P(B}
{A}
{$B8E$$%5!<%P$O%G%#%9%/>c32Ey$,$"$C$?$N$G!"<wL?$HH=CG$5$l!"%G!<%?$,B;<:$9$kA0$K?7$7$$%5!<%P$K0\9T$5$l$^$7$?!#(Bftp-master.debian.org$B$O(BDebian$B$N(B official package $B%j%]%8%H%j$G$9!#%Q%C%1!<%8$N%9%]%s%5!<(B(mentor)$B$rC5$9$N$O(Bmentors.debian.net$B!#(B }

\santaku
{pbuilder$B$K(Bclang support$B$,DI2C$5$l$^$7$?!#C/$,=q$$$?%Q%C%A$G$7$g$&$+!)(B}
{Sylvestre Ledru}
{Junichi Uekawa}
{Hideki Yamane}
{C}
{Debian$B$N(BClang$B%5%]!<%H$OCe!9$H?J$s$G$$$^$9!#(B}

\santaku
{DPN - 2013$BG/(B3$B7n(B4$BF|9f$K<h$j>e$2$i$l$?F|K\$N%$%Y%s%H$O(B}
{Open Source Conference 2013 Tokyo/Spring}
{Open Source Conference 2013 Hamamatu}
{Open Source Conference 2013 Tokushima}
{A}
{\url{http://henrich-on-debian.blogspot.jp/2013/02/open-source-conference-2013-tokyospring.html} $B>\:Y$O8e$[$I!#(B}


%\end{multicols}


% % (query-replace-regexp "<.*?>" "")
% % (query-replace-regexp "^[    ]\+" "")

%-------------------------------------------------------------------------------
\dancersection{git-buildpackage を使ってみる}{杉本典充}
%-------------------------------------------------------------------------------

\subsection{はじめに}

debian パッケージを作成し管理する方法はいろいろあると思います。

今回はgitでパッケージを管理し、ビルドする仕組みを提供する git-buildpackage (gbp) に
ついて調べてみました。

\subsection{git-buildpackage とは}

debian では「git-buildpackage」パッケージを提供しています。
sid では git-buildpackage-0.9.25 を提供しており、gbp コマンド\footnote{中身は python3 で書かれています。}、
git-pbuilder コマンドが含まれています。

gbp コマンドには debian パッケージを開発する上で様々な操作ができる便利な以下のサブコマンドがあります\footnote{\$ gbp --list-cmds}。

\begin{itemize}
  \item	gbp buildpackage
  \item gbp import-orig
  \item gbp export-orig
  \item gbp import-{dsc,dscs}
  \item gbp import-ref
  \item gbp dch
  \item gbp pq
  \item gbp {pull,clone,push}
  \item gbp create-remote-repo
  \item gbp tag
  \item gbp config
  \item gbp pristine-tar
  \item gbp setup-gitattributes
\end{itemize}

なお、git-buildpackage というパッケージ名になっていますが、他のバージョン管理システムに対応した cvs-buildpackage、svn-buildpackage、mercurial-buildpackage パッケージも存在します。


\subsection{git-buildpackage を使って tarball を新規にパッケージ化する}

ここでは例として、GNU hello\footnote{\url{https://www.gnu.org/software/hello/}} の tarball を debian パッケージ化\footnote{\url{https://packages.debian.org/ja/sid/hello} が既に存在します}する作業を gbp コマンドを使って行ってみます。

なお、今回の作業環境は Debian sid とします。

\subsubsection{必要な debian パッケージのインストール}

まずは git-buildpackage パッケージを apt-get コマンドでインストールします。
そして、debian パッケージの作成に必要な build-essential、大変便利な debhelper をインストールします。

\begin{commandline}
$ sudo apt-get update
$ sudo apt-get install git-buildpackage
$ sudo apt-get install build-essential debhelper dh-make
\end{commandline}

今回 Debian パッケージを作成する GNU hello のビルドに使うツール、およびビルドに必要な依存パッケージをインストールしておきます。

\begin{commandline}
$ sudo apt-get install wget
$ sudo apt-get install gnulib help2man
\end{commandline}


\subsubsection{git リポジトリの作成と tarball のインポート}

まずはGNU hello の tarball をダウンロードします。

\begin{commandline}
$ wget https://ftp.gnu.org/gnu/hello/hello-2.10.tar.gz
\end{commandline}

git のローカルブランチのディレクトリ名は tarball のファイル名と同じ "hello" にすることとし、git リポジトリ "hello" をローカルに作成します。最近の git の使い方としてデフォルトを main ブランチにするように変わってきていますので、ブランチ名を指定しています\footnote{git-buildpackage ではデフォルトのブランチ名は現状 master のままになっています。}。

\begin{commandline}
$ git init hello --initial-branch=main
$ cd hello
\end{commandline}

tarball を指定して gbp import-orig コマンドを実行すると、tarball を git に取り込んでコミットします。デフォルトでは master ブランチを作成しますが、\url{https://salsa.debian.org/} の git デフォルトブランチは main のため --debian-branch=main を指定しています。

\begin{commandline}
$ gbp import-orig --debian-branch=main --pristine-tar ../hello-2.10.tar.gz

What will be the source package name? [hello]
What is the upstream version? [2.10]
gbp:info: Importing '../hello-2.10.tar.gz' to branch 'upstream'...
gbp:info: Source package is hello
gbp:info: Upstream version is 2.10
gbp:info: Successfully imported version 2.10 of ../hello_2.10.orig.tar.gz
\end{commandline}

gbp import-orig で tarball を取り込むと以下のブランチとタグが作成されます。

\begin{commandline}
$ git branch
* main
  pristine-tar
  upstream
\end{commandline}

\begin{commandline}
$ git tag
upstream/2.10
\end{commandline}

git-buildpackage において git のブランチは以下の使い分けをします。

\begin{itemize}
  \item master (または main)ブランチ：debian パッケージ全体の内容を含むブランチ
  \item upstreamブランチ：upstream の配布物をそのまま展開した内容を含むブランチ
  \item pristine-tarブランチ：debian パッケージのビルドに使う upstream ブランチから生成する orig.tar.gz と upstream の tarball の差分をバイナリで保持するブランチ
\end{itemize}

gbp import-orig コマンドに --pristine-tar オプションを指定した場合は pristine-tar ブランチも作成されます。

pristine-tar ブランチの作成は必須ではありません。ただ、tarball を展開して upstream ブランチに保存して、逆に upstream ブランチのファイルから再生成した tarball ファイルのハッシュ値が、gbp import-orig で取り込んだ tarball のハッシュ値と同じにならないことがあります。
pristine-tar ブランチにはその差分をバイナリで保存して、gbp buildpackage で debian パッケージをビルドするときに適用することでハッシュ値が同じになるように調整しています。

\begin{commandline}
$ git checkout pristine-tar
Switched to branch 'pristine-tar'
Your branch is up to date with 'origin/pristine-tar'.

$ ls -l
合計 28
-rw-r--r-- 1 norimitu norimitu 6107  3月 19 04:25 hello_2.10.orig.tar.gz.delta
-rw-r--r-- 1 norimitu norimitu   41  3月 19 04:25 hello_2.10.orig.tar.gz.id

$ cat hello_2.11.orig.tar.gz.id
61732d00c5e503e9ced7f7f57f5b1c14a449c95c

$ file hello_2.11.orig.tar.gz.delta
hello_2.11.orig.tar.gz.delta: gzip compressed data, from Unix, original size modulo 2^32 30720

$ git checkout master
\end{commandline}


\subsubsection{debian ディレクトリを生成する}

debian パッケージを作成するのは、main ブランチに debian ディレクトリを作成します。
dh\_make コマンドを使って作成できます。

\begin{commandline}
$ dh_make -p hello_2.10

Type of package: (single, indep, library, python)
[s/i/l/p]?
Maintainer Name     : Norimitsu Sugimoto
Email-Address       : dictoss@live.jp
Date                : Thu, 17 Mar 2022 13:27:56 +0000
Package Name        : hello
Version             : 2.10
License             : blank
Package Type        : single
Are the details correct? [Y/n/q]
Skipping creating ../hello_2.10.orig.tar.gz because it already exists
Done. Please edit the files in the debian/ subdirectory now.
\end{commandline}

dch コマンドで debian/changelog ファイルを開いて編集します\footnote{通常、debian パッケージを新規作成する場合は ITP のバグ報告を行っている前提があり、debian/changelog ファイルにそのバグ番号を含めます。}。

\begin{commandline}
$ dch

hello (2.10-1) UNRELEASED; urgency=medium

  * Initial release (Closes: #nnnn)  <nnnn is the bug number of your ITP>

 -- Norimitsu Sugimoto <dictoss@live.jp>  Thu, 17 Mar 2022 13:31:04 +0000
\end{commandline}

debian/control ファイルを編集し、パッケージのメタ情報を更新します。

\begin{commandline}
$ vim debian/control

Source: hello
Section: devel
Priority: optional
Maintainer: Norimitsu Sugimoto <dictoss@live.jp>
Build-Depends: debhelper-compat (= 13), autotools-dev, gnulib, help2man
Standards-Version: 4.6.0
Homepage: http://www.gnu.org/software/hello/
#Vcs-Browser: https://salsa.debian.org/debian/hello
#Vcs-Git: https://salsa.debian.org/debian/hello.git
Rules-Requires-Root: no

Package: hello
Architecture: any
Depends: ${shlibs:Depends}, ${misc:Depends}
Description: example package based on GNU hello
 The GNU hello program produces a familiar, friendly greeting.  It
 allows non-programmers to use a classic computer science tool which
 would otherwise be unavailable to them.
 .
 Seriously, though: this is an example of how to do a Debian package.
 It is the Debian version of the GNU Project's `hello world' program
 (which is itself an example for the GNU Project).
\end{commandline}

debian/rules ファイルを編集し、debhelper を使って debian パッケージをビルドできるように調整します。

\begin{commandline}
$ vim debian/rules

%:
        dh $@

override_dh_auto_clean:
        [ ! -f Makefile ] || $(MAKE) distclean
\end{commandline}

それでは gbp buildpackage コマンドを使って debian パッケージをビルドしてみます。しかし、コミットしていないファイルがあるとコマンドがエラーになります。まずは編集した debian ディレクトリを git へコミットします。

\begin{commandline}
$ gbp buildpackage --git-pristine-tar --git-debian-branch=main

gbp:error: You have uncommitted changes in your source tree:
gbp:error: On branch master
Untracked files:
  (use "git add <file>..." to include in what will be committed)
        debian/

nothing added to commit but untracked files present (use "git add" to track)

gbp:error: Use --git-ignore-new to ignore.
\end{commandline}

\begin{commandline}
$ git add debian
$ git commit -m "add debian directory."
\end{commandline}

再度、gbp buildpackage コマンドを実行します。debian パッケージのビルドに成功しました。しかし、lintian の警告がたくさん出ていますのでリリースするまでに警告をできるだけ減らすパッケージを見直す必要があります。

\begin{commandline}
$ gbp buildpackage --git-pristine-tar --git-debian-branch=main

gbp:info: Performing the build
dpkg-buildpackage -us -uc -ui -i -I
dpkg-buildpackage: info: source package hello
dpkg-buildpackage: info: source version 2.10-1
dpkg-buildpackage: info: source distribution UNRELEASED
dpkg-buildpackage: info: source changed by Norimitsu Sugimoto <dictoss@live.jp>
 dpkg-source -i -I --before-build .
dpkg-buildpackage: info: host architecture amd64
 debian/rules clean

  (省略)

dpkg-buildpackage: info: full upload (original source is included)
Now running lintian hello_2.10-1_amd64.changes ...
E: hello: changelog-is-dh_make-template
  (省略)
W: hello: wrong-bug-number-in-closes #nnnn in the installed changelog (line 4)
Finished running lintian.
\end{commandline}

これで debian パッケージのビルドができました。
ビルド時に生成されたファイルを削除してから、main ブランチにタグを作成し、
リモートの git リポジトリに push します。

\begin{commandline}
$ gbp buildpackage --git-tag-only --git-debian-branch=main
gbp:info: Tagging Debian package 2.10-1 as debian/2.10-1 in git

$ git tag
debian/2.10-1
upstream/2.10
\end{commandline}

\begin{commandline}
$ git remote add origin git@salsa.debian.org:dictoss-guest/hello.git
$ git remote -v
origin  git@salsa.debian.org:dictoss-guest/hello.git (fetch)
origin  git@salsa.debian.org:dictoss-guest/hello.git (push)

$ git push -u origin --all
$ git push -u origin --tag
\end{commandline}


% 過去のパッケージの再ビルド手順

\subsection{upstreamで新しいバージョンのtarballがリリースされた場合の作業}

まず、更新された tarball をダウンロードします。

\begin{commandline}
$ cd ..
$ wget https://ftp.gnu.org/gnu/hello/hello-2.11.tar.gz
\end{commandline}

gbp import-orig コマンドで tarball を git のローカルブランチに取り込みます。

\begin{commandline}
$ gbp import-orig --debian-branch=main --pristine-tar ../hello-2.11.tar.gz

What is the upstream version? [2.11]
gbp:info: Importing '../hello-2.11.tar.gz' to branch 'upstream'...
gbp:info: Source package is hello
gbp:info: Upstream version is 2.11
gbp:info: Replacing upstream source on 'main'
gbp:info: Successfully imported version 2.11 of ../hello_2.11.orig.tar.gz
\end{commandline}

git branch を実行するとブランチの数に変化はありません。

\begin{commandline}
$ git branch
* main
  pristine-tar
  upstream
\end{commandline}

git tag を実行すると "upstream/2.11" が追加されています。

\begin{commandline}
$ git tag

debian/2.10-1
upstream/2.10
upstream/2.11
\end{commandline}

dch コマンドで debian/changelog ファイルを開いて編集します。

\begin{commandline}
$ dch

hello (2.11-1) UNRELEASED; urgency=medium

  * New upstream version 2.11

 -- Norimitsu Sugimoto <dictoss@live.jp>  Fri, 19 Mar 2022 14:10:38 +0000

hello (2.10-1) UNRELEASED; urgency=medium

  * Initial release (Closes: #nnnn)  <nnnn is the bug number of your ITP>

 -- Norimitsu Sugimoto <dictoss@live.jp>  Fri, 18 Mar 2022 13:39:37 +0000
\end{commandline}

debian ディレクトリの変更内容をコミットします。

\begin{commandline}
$ git add debian
$ git commit -m "New upstream version 2.11."
\end{commandline}

gbp buildpackage を実行して debian パッケージをビルドしてみるとエラーになりました。

\begin{commandline}
$ gbp buildpackage --git-pristine-tar --git-debian-branch=main

  （省略）
*** error: gettext infrastructure mismatch: using a Makefile.in.in from gettext
version 0.18 but the autoconf macros are from gettext version 0.20
make[3]: *** [Makefile:165: check-macro-version] Error 1
make[3]: *** Waiting for unfinished jobs....
（省略）
dh_auto_build: error: make -j3 returned exit code 2
make: *** [debian/rules:18: binary] エラー 25
dpkg-buildpackage: error: debian/rules binary subprocess returned exit status 2
debuild: fatal error at line 1182:
dpkg-buildpackage -us -uc -ui -i -I failed
gbp:error: 'debuild -i -I' failed: it exited with 29
\end{commandline}

gettext のバージョンが sid のバージョンと tarball のファイルに書いてあるバージョンが
合わないためエラーになったようです。gbp pq コマンドを使ってパッチ作成用のブランチを作ります。
gbp pq import コマンドを実行すると patch-queue/main ブランチが作成されチェックアウトした状態になります。

\begin{commandline}
$ gbp pq import
gbp:info: Trying to apply patches at 'cacbfce059a86c836513ff70eb6af780d9255910'
gbp:info: 0 patches listed in 'debian/patches/series' imported on 'patch-queue/main'

$ git branch
  main
* patch-queue/main
  pristine-tar
  upstream
\end{commandline}

patch-queue/main ブランチ上で変更が必要なファイルを編集します。

\begin{commandline}
$ vim configure.ac
$ git diff
diff --git a/configure.ac b/configure.ac
index 6d52c41..f101d7b 100644
--- a/configure.ac
+++ b/configure.ac
@@ -62,7 +62,7 @@ dnl Ensure VLAs are not used.
 AC_DEFINE([GNULIB_NO_VLA], [1], [Define to 1 to disable use of VLAs])

 dnl i18n support from GNU gettext.
-AM_GNU_GETTEXT_VERSION([0.18.1])
+AM_GNU_GETTEXT_VERSION([0.21])
 AM_GNU_GETTEXT([external])

 AC_CONFIG_FILES([Makefile
\end{commandline}

patch-queue/main ブランチに変更したファイルをコミットします。

\begin{commandline}
$ git add configure.ac
$ git commit -m "change to gettext version to use with sid."
[patch-queue/main cc568aa] change to gettext version to use with sid.
 1 file changed, 1 insertion(+), 1 deletion(-)
\end{commandline}

すべてのファイルの変更が完了した後、gbp pq export コマンドを実行して編集を終了します。
gbp pq export コマンドを実行すると main ブランチをチェックアウトした状態になり、
main ブランチと patch-queue/main ブランチの差分が debian/patches/ ディレクトリに
生成されます。

\begin{commandline}
$ gbp pq export
gbp:info: On 'patch-queue/main', switching to 'main'
gbp:info: Generating patches from git (main..patch-queue/main)
\end{commandline}

\begin{commandline}
$ git branch
* main
  patch-queue/main
  pristine-tar
  upstream
\end{commandline}

\begin{commandline}
$ git status .
On branch main
Your branch is ahead of 'origin/main' by 3 commits.
  (use "git push" to publish your local commits)

Untracked files:
  (use "git add <file>..." to include in what will be committed)
        debian/patches/

nothing added to commit but untracked files present (use "git add" to track)

$ ls -l debian/patches/
合計 8
-rw-r--r-- 1 norimitu norimitu 596  3月 19 02:28 0001-change-to-gettext-version-to-use-with-sid.patch
-rw-r--r-- 1 norimitu norimitu  53  3月 19 02:28 series
\end{commandline}

生成された debian/patches ディレクトリを main ブランチにコミットします。

\begin{commandline}
$ git add debian/patches
$ git commit -m "change to gettext version to use with sid."
[main 6ae3e38] change to gettext version to use with sid.
 2 files changed, 22 insertions(+)
 create mode 100644 debian/patches/0001-change-to-gettext-version-to-use-with-sid.patch
 create mode 100644 debian/patches/series
\end{commandline}

gbp buildpackage を実行すると、更新した debian パッケージをビルドできます。

\begin{commandline}
$ gbp buildpackage --git-pristine-tar --git-debian-branch=main

$ ls ..
hello                          hello_2.10-1.debian.tar.xz    hello_2.10-1_amd64.deb      hello_2.11-1_amd64.buildinfo
hello-2.10.tar.gz              hello_2.10-1.dsc              hello_2.10.orig.tar.gz      hello_2.11-1_amd64.changes
hello-2.11.tar.gz              hello_2.10-1_amd64.build      hello_2.11-1.debian.tar.xz  hello_2.11-1_amd64.deb
hello-dbgsym_2.10-1_amd64.deb  hello_2.10-1_amd64.buildinfo  hello_2.11-1.dsc            hello_2.11.orig.tar.gz
hello-dbgsym_2.11-1_amd64.deb  hello_2.10-1_amd64.changes    hello_2.11-1_amd64.build
\end{commandline}

ビルド時に生成されたファイルを削除してから、main ブランチにタグを作成し、
リモートの git リポジトリに push します。

\begin{commandline}
$ gbp buildpackage --git-tag-only --git-debian-branch=main
gbp:info: Tagging Debian package 2.11-1 as debian/2.11-1 in git

$ git tag
debian/2.10-1
debian/2.11-1
upstream/2.10
upstream/2.11
\end{commandline}

\begin{commandline}
$ git push -u origin main upstream pristine-tar
$ git push -u origin --tag
\end{commandline}

作業を終えた patch-queue/main ブランチは不要のため削除します。

\begin{commandline}
$ git branch -D patch-queue/main
\end{commandline}


\subsection{今回取り扱えなかったこと}

今回取り扱えなかったことはたくさんあります。

\begin{itemize}
  \item upstream が git リポジトリのみで tarball を提供していない場合
  \item sid、stable、stable-backports、experimental を平行で管理する方法
  \item gbp buildpackage --git-pbuilder
  \item gbp.conf
  \item upstream のソースコードにすでに debian ディレクトリが存在している場合
\end{itemize}


\subsection{おわりに}

git-buildpackage パッケージに入っている gbp コマンドを調べてみました。
quilt コマンドを使ったパッチの作成を gbp pq コマンドでより便利に使えるのがよいと思いました。

実際にパッケージをメンテナンスするときには、複数のリリースを 1 つの git リポジトリで管理したいのだと思います。その場合にどのように git-buildpackage を使いこなせばよいか、まだまだ勉強が必要であると感じました。


\subsection{参考文献}

\begin{itemize}
  \item Debian Wiki - PackagingWithGit \url{https://wiki.debian.org/PackagingWithGit}
  \item Debian パッケージング道場 git-buildpackage\\ \url{https://tokyodebian-team.pages.debian.net/pdf2015/debianmeetingresume201509-presentation.pdf}
  \item Debian 新メンテナーガイド 第6章 パッケージのビルド \\ \url{https://www.debian.org/doc/manuals/maint-guide/build.ja.html}
  \item git-buildpackageを用いたdebパッケージ管理方法の紹介 \\ \url{https://blog.cybozu.io/entry/2019/04/11/110000}
\end{itemize}

% 冊子にするために、4の倍数にする必要がある。
% そのための調整
\dancersection{メモ}{}
\mbox{}\newpage
\mbox{}\newpage
\mbox{}\newpage

\vspace*{15cm}
\hrule
\vspace{2mm}
\includegraphics[width=2cm]{image-assets/openlogo-nd.eps}
\noindent \Large \bf Debian 勉強会資料\\
\noindent \normalfont \debmtgyear{}年\debmtgmonth{}月\debmtgdate{}日 \hspace{5mm}  初版第1刷発行\\
\noindent \normalfont 東京エリア Debian 勉強会 （編集・印刷・発行）\\
\hrule
\end{document}
